\ECchapter{01}{Renewable Energy}{Can solar power meet tomorrow's energy needs?}{ECOrange}

\ECObjectives{
\begin{itemize}[leftmargin=*,itemsep=1pt,topsep=1pt]
\item Explain how a photovoltaic cell and a grid-connected PV system operate.
\item Interpret capacity data and estimate annual production.
\item Discuss intermittency, storage, grid integration and recycling.
\end{itemize}
}

\begin{multicols}{2}

\begin{engineeringnews}
Global renewable capacity grew rapidly in 2025, and solar power represented most of the new installations. The challenge is no longer limited to producing more panels: engineers must also connect variable generation to stronger, smarter and more flexible electricity grids.
\end{engineeringnews}

\begin{ecarticle}
A photovoltaic (PV) cell converts light directly into electricity. Most commercial cells are made from silicon, a semiconductor. When photons transfer enough energy to the material, electrons can move. The internal electric field of the p--n junction separates positive and negative charges and creates a direct current.

A PV module connects many cells together. Modules are then combined into strings and arrays to reach the required voltage and power. However, a module does not operate at a fixed point. Its voltage-current characteristic changes with solar irradiance and temperature. A maximum power point tracker adjusts the operating point so that the array can deliver as much power as possible.

The generated electricity is direct current, whereas buildings and public grids use alternating current. An inverter therefore converts DC into AC, synchronises voltage and frequency with the grid, and provides protection and monitoring functions. A battery can store surplus electricity, but it also increases cost, material use and conversion losses.

Solar power is variable because output changes with clouds, seasons and time of day. Engineers integrate large amounts of PV through forecasting, flexible demand, interconnections, storage and controllable generation. The objective is not to make every building independent, but to balance the complete electricity system.

The environmental impact of a PV system must be assessed over its whole life cycle. Manufacturing requires energy and materials, while end-of-life modules contain glass, aluminium, copper, polymers and semiconductor materials. Designing panels that are efficient, durable, repairable and recyclable is therefore a major engineering objective.
\end{ecarticle}

\begin{figure}[H]
\centering
\includegraphics[width=\linewidth]{images/EC01/topaz_solar_farm.jpg}
\caption*{\footnotesize Topaz Solar Farm, California. NASA image, public domain.}
\end{figure}

\columnbreak

\begin{engineeringfocus}
\textbf{Designing a photovoltaic installation}

Engineers must optimise several parameters at the same time:
\begin{itemize}[leftmargin=*,itemsep=1pt]
\item panel orientation and tilt;
\item inverter sizing and MPPT architecture;
\item cable cross-sections and electrical losses;
\item thermal behaviour and ventilation;
\item protection and isolation devices;
\item monitoring and maintenance strategy;
\item expected lifetime and recyclability.
\end{itemize}

A good design must combine safety, reliability, high annual energy yield and acceptable cost.
\end{engineeringfocus}

\begin{figure}[H]
\centering
\includegraphics[width=\linewidth]{images/EC01/pv_system.png}
\caption*{\footnotesize Simplified residential PV system. S-kei, CC0.}
\end{figure}

\begin{tcolorbox}[colback=white,colframe=ECBlue,title=\sffamily\bfseries Global solar PV capacity]
\centering
\begin{tikzpicture}
\begin{axis}[
width=.96\linewidth,height=45mm,
xlabel={Year},ylabel={Installed capacity (GW)},
xmin=2016,xmax=2025,ymin=0,ymax=2500,
xtick={2016,2018,2020,2022,2024},
grid=major,
tick label style={font=\scriptsize},
label style={font=\scriptsize}]
\addplot+[mark=*,thick] table[x=year,y=capacity_gw,col sep=comma]{data/EC01/global_solar_pv_capacity.csv};
\end{axis}
\end{tikzpicture}

{\scriptsize Source: IRENA, Renewable Capacity Statistics 2026.}
\end{tcolorbox}

\begin{tcolorbox}[colback=yellow!10,colframe=ECOrange,title=\sffamily\bfseries Engineering Insight]
The Earth receives an enormous amount of solar energy, but useful production depends on available area, local weather, orientation, conversion efficiency and the ability of the grid to absorb variable power.
\end{tcolorbox}

\begin{tcolorbox}[colback=ECLightBlue,colframe=ECBlue,title=\sffamily\bfseries Key Figures]
\begin{tabularx}{\linewidth}{@{}Xr@{}}
Commercial module efficiency & 20--24\%\\
Typical module lifetime & 25--30 years\\
Typical inverter efficiency & 97--99\%\\
Annual module degradation & 0.3--0.5\%/year\\
Global PV capacity in 2025 & $\approx 2.4$ TW
\end{tabularx}
\end{tcolorbox}

\begin{vocabulary}
\begin{tabularx}{\linewidth}{@{}lX@{}}
solar cell & cellule photovoltaïque\\
irradiance & irradiance\\
inverter & onduleur\\
grid & réseau électrique\\
storage & stockage\\
yield & production énergétique\\
lifetime & durée de vie\\
recycling & recyclage
\end{tabularx}
\end{vocabulary}

\begin{questions}
\item Explain the role of the p--n junction.
\item Why does a PV installation need an inverter?
\item What is the purpose of MPPT control?
\item Use the graph to estimate the increase between 2020 and 2025.
\item Give three solutions for integrating variable solar power.
\item Why must environmental assessment include manufacturing and end-of-life?
\end{questions}

\begin{challenge}
Your school consumes 410 MWh of electricity each year and has 850 m$^2$ of usable roof area. Prepare a preliminary PV study: estimate the peak power, estimate annual production, discuss self-consumption, decide whether a battery is justified and propose useful monitoring indicators.
\end{challenge}

\begin{applicationexercise}
A 420~W PV module receives an average irradiance of 850~W/m$^2$ over an active area of 2.0~m$^2$. Calculate its instantaneous efficiency. Then estimate the electrical energy produced by 120 identical modules during 4.5 equivalent full-sun hours. State two reasons why the real daily yield may be lower.
\end{applicationexercise}

\begin{speaking}
Working in pairs, prepare a five-minute presentation answering the question: ``Can Europe rely mainly on solar electricity by 2050?'' Support your arguments with technical evidence, environmental impacts and economic considerations.
\end{speaking}

\end{multicols}